% =============================================================================
% APPENDIX BL: DOMAIN-CRISIS: Geopolitical Crisis Escalation Modeling
% =============================================================================
% Category: DOMAIN
% Module: BCM2_07_GEO (Geopolitical Domain)
% Version: 1.0 (January 2026)
% Status: Draft
% Language: English
%
% This appendix provides the formal framework for modeling geopolitical crisis
% escalation using sequential game theory and behavioral economics. It documents
% the methodology behind MOD-GEO-001 (US-Iran Crisis Escalation Model).
%
% =============================================================================

% =============================================================================
% CROSS-REFERENCE STRUCTURE
% =============================================================================
%
% CHAPTER LINKAGE:
% - Primary Chapter: Chapter 14 (Policy Applications)
% - Secondary Chapters: Chapter 9 (Context), Chapter 5 (Complementarity)
%
% DEPENDENCIES (what this appendix REQUIRES):
% - Appendix UNMAPPED_B (CORE-HOW): Complementarity framework, γ-matrix
% - Appendix UNMAPPED_V (CORE-WHEN): Context dimensions Ψ
% - Appendix UNMAPPED_MS (REF-MODELSPACE): Theory catalog (Game Theory)
% - Appendix UNMAPPED_BBB (CORE-WHERE): Parameter sources
%
% DEPENDENTS (what REQUIRES this appendix):
% - MOD-GEO-001: US-Iran Crisis Model
% - Future geopolitical models (Taiwan, Korea, NATO)
%
% =============================================================================

\chapter{Geopolitical Crisis Escalation Modeling}
\label{app:domain-crisis}

% -----------------------------------------------------------------------------
% CHAPTER LINKAGE BOX
% -----------------------------------------------------------------------------

\begin{tcolorbox}[colback=purple!5!white, colframe=purple!75!black,
    title=Chapter Linkage]

\textbf{Primary Chapter:} Chapter 14: Policy Applications
\begin{itemize}[nosep]
\item Section 14.3: Crisis Decision-Making $\leftrightarrow$ This Appendix Section \ref{sec:crisis-framework}
\item Section 14.5: International Relations $\leftrightarrow$ This Appendix Section \ref{sec:crisis-game}
\end{itemize}

\textbf{Secondary Chapters:}
\begin{itemize}[nosep]
\item Chapter 9: Context (Ψ-Dimensions) --- provides context framework
\item Chapter 5: Complementarity (γ) --- provides interaction structure
\end{itemize}

\textbf{Reading Guide:}
\begin{itemize}[nosep]
\item For conceptual overview: Read Chapter 14 first
\item For game-theoretic details: This appendix provides complete specification
\item For parameter sources: See Appendix UNMAPPED_BBB (CORE-WHERE)
\end{itemize}

\end{tcolorbox}

% -----------------------------------------------------------------------------
% ABSTRACT
% -----------------------------------------------------------------------------

\begin{abstract}
This appendix develops a formal framework for modeling geopolitical crisis escalation using sequential game theory integrated with the EBF context framework. We specify how state actors make decisions under incomplete information, how context dimensions ($\Psi$) modify utility functions and beliefs, and how to derive testable predictions using Subgame Perfect Equilibrium (SPE). The framework is demonstrated through the US-Iran Crisis Model (MOD-GEO-001), showing how behavioral economics enriches traditional international relations theory.
\end{abstract}

% -----------------------------------------------------------------------------
% QUICK REFERENCE BOX
% -----------------------------------------------------------------------------

\begin{tcolorbox}[colback=blue!5!white, colframe=blue!75!black,
    title=Quick Reference: Crisis Escalation Framework]

\textbf{Core Equation:}
\begin{equation}
U^*_i(\text{outcome}) = U_i(\text{outcome}) \cdot \prod_k f_k(\Psi_k)
\end{equation}

\textbf{Key Components:}
\begin{itemize}[nosep]
\item \textbf{Players:} State actors with utility functions $U_i$
\item \textbf{Strategies:} Action sets $S_i = \{s_1, s_2, ..., s_n\}$
\item \textbf{Context:} Ψ-modifiers adjust base utilities
\item \textbf{Solution:} Subgame Perfect Equilibrium via backward induction
\end{itemize}

\textbf{Model Template:} MOD-GEO-001 (US-Iran Crisis)

\textbf{Prediction Output:} $P(\text{action} | \text{context}, \text{trigger})$

\end{tcolorbox}

% -----------------------------------------------------------------------------
% APPENDIX SCOPE BOX
% -----------------------------------------------------------------------------

\begin{tcolorbox}[colback=gray!5!white, colframe=gray!75!black,
    title=Appendix Scope]

\textbf{Ziel:} Formales Framework für die Modellierung geopolitischer Krisen-Eskalation mit verhaltensökonomischer Integration.

\textbf{In-Scope:}
\begin{itemize}[nosep]
\item Sequential Game Formalisierung für Krisenentscheidungen
\item Integration von Ψ-Kontext-Dimensionen in Utility-Funktionen
\item Parametrisierung via LLMMC + Bayesian Update
\item Ableitung testbarer Vorhersagen mit Konfidenzintervallen
\end{itemize}

\textbf{Out-of-Scope:}
\begin{itemize}[nosep]
\item Militärische Operationsplanung
\item Detaillierte Länderprofile (siehe BCM2\_07\_GEO)
\item Historische Fallstudien (siehe Case Registry)
\end{itemize}

\textbf{Constraints:}
\begin{itemize}[nosep]
\item Parameter-Unsicherheit: ±15-25\% auf allen Schätzungen
\item Modell gilt für bilaterale Krisen (Multi-Actor erfordert Erweiterung)
\item Validierung erst nach Outcome möglich
\end{itemize}

\end{tcolorbox}


% =============================================================================
% SECTION 1: THE FUNDAMENTAL QUESTION
% =============================================================================

\section{The Fundamental Question}
\label{sec:crisis-question}

\begin{tcolorbox}[colback=yellow!5!white, colframe=yellow!75!black,
    title=Central Question]
\centering
\textbf{Wie können wir die Wahrscheinlichkeit von Eskalation in geopolitischen Krisen quantifizieren, unter Berücksichtigung von Kontext, kognitiven Verzerrungen und strategischer Interaktion?}
\end{tcolorbox}

Traditionelle Ansätze der Internationalen Beziehungen (Realismus, Liberalismus, Konstruktivismus) bieten qualitative Frameworks, aber wenig präzise Vorhersagen. Das EBF Crisis Framework integriert:

\begin{enumerate}
\item \textbf{Game Theory:} Strategische Interaktion mit Gleichgewichtskonzepten
\item \textbf{Behavioral Economics:} Kognitive Verzerrungen, Loss Aversion, Face-Saving
\item \textbf{Context Framework:} Ψ-Dimensionen als systematische Modifikatoren
\item \textbf{Bayesian Updating:} LLMMC-Priors + Evidenz = Posterior-Vorhersagen
\end{enumerate}


% =============================================================================
% SECTION 2: THE CRISIS ESCALATION FRAMEWORK
% =============================================================================

\section{The Crisis Escalation Framework}
\label{sec:crisis-framework}

\subsection{Axioms}

\begin{axiom}[Rational Actor Approximation]
\label{ax:rational}
State actors maximize expected utility given beliefs about other actors' strategies, subject to cognitive and contextual constraints.
\end{axiom}

\begin{axiom}[Context Dependency]
\label{ax:context}
Base utility functions are modified by context dimensions $\Psi = \{\Psi_T, \Psi_S, \Psi_E, \Psi_I, \Psi_K, \Psi_C, \Psi_M, \Psi_F\}$.
\end{axiom}

\begin{axiom}[Sequential Rationality]
\label{ax:sequential}
At each decision node, actors choose strategies that maximize expected utility given beliefs about future play (Subgame Perfect Equilibrium).
\end{axiom}

\begin{axiom}[Incomplete Information]
\label{ax:incomplete}
Actors have beliefs about other actors' types (e.g., «Hawk» vs. «Dove») that influence equilibrium play.
\end{axiom}

\subsection{Formal Definition}

A \textbf{Crisis Escalation Game} is defined as a tuple:

\begin{equation}
\Gamma = \langle N, H, P, (U_i)_{i \in N}, (\Psi_k)_{k \in K}, (f_k)_{k \in K}, \mu \rangle
\end{equation}

Where:
\begin{itemize}
\item $N = \{1, 2, ..., n\}$: Set of players (state actors)
\item $H$: Set of histories (sequences of actions)
\item $P: H \rightarrow N$: Player function (who moves at each history)
\item $U_i: Z \rightarrow \mathbb{R}$: Utility function for player $i$ over terminal nodes $Z$
\item $\Psi_k$: Context dimension $k$ (temporal, social, economic, etc.)
\item $f_k: \Psi_k \rightarrow \mathbb{R}^+$: Context modifier function
\item $\mu$: Belief system over information sets
\end{itemize}


% =============================================================================
% SECTION 3: UTILITY FUNCTIONS WITH CONTEXT MODIFIERS
% =============================================================================

\section{Utility Functions with Context Modifiers}
\label{sec:crisis-utility}

\subsection{Base Utility Structure}

For a state actor $i$, the base utility function is:

\begin{equation}
U_i(\text{outcome}) = \sum_j \alpha_{ij} \cdot \text{Benefit}_j(\text{outcome}) - \sum_k \beta_{ik} \cdot \text{Cost}_k(\text{outcome})
\end{equation}

\textbf{Typical Benefit Dimensions:}
\begin{itemize}
\item Credibility / Reputation
\item Security / Threat Reduction
\item Domestic Political Win
\item Economic Gain
\end{itemize}

\textbf{Typical Cost Dimensions:}
\begin{itemize}
\item War Costs (casualties, resources)
\item Economic Disruption
\item Escalation Risk
\item International Isolation
\end{itemize}

\subsection{Context-Modified Utility}

The effective utility incorporates context modifiers:

\begin{equation}
U^*_i(\text{outcome}) = U_i(\text{outcome}) \cdot \prod_k f_k(\Psi_k)
\end{equation}

\textbf{Context Modifier Functions:}

\begin{align}
f(\Psi_T) &= 1 + \gamma_T \cdot (\text{temporal\_pressure}) \\
f(\Psi_S) &= 1 - \gamma_S \cdot (1 - \text{allied\_support}) \\
f(\Psi_E) &= 1 + \gamma_E \cdot (\text{resource\_readiness}) \\
f(\Psi_K) &= 1 + \gamma_K \cdot (\text{face\_saving\_necessity}) \\
f(\Psi_C) &= 1 + \gamma_C \cdot (\text{leader\_impulsivity})
\end{align}

Where $\gamma_k$ are sensitivity parameters estimated via LLMMC.


% =============================================================================
% SECTION 4: THE SEQUENTIAL GAME STRUCTURE
% =============================================================================

\section{The Sequential Game Structure}
\label{sec:crisis-game}

\subsection{Extensive Form Representation}

A typical bilateral crisis game has the following structure:

\begin{verbatim}
                         Actor 1
                        /   |   \
                      /     |     \
                THREATEN  STRIKE  NEGOTIATE
                   |        |         |
                Actor 2   Actor 2   Actor 2
               /  |  \    /  |  \   /  |  \
             ...  ... ... ... ... ... ... ... ...
\end{verbatim}

\subsection{Strategy Sets}

\textbf{Typical Aggressor Strategies:}
\begin{itemize}
\item $s_1$: Status Quo (maintain pressure)
\item $s_2$: Threaten (escalate rhetoric)
\item $s_3$: Limited Strike (targeted action)
\item $s_4$: Major Strike (significant action)
\item $s_5$: Negotiate (seek deal)
\end{itemize}

\textbf{Typical Defender Strategies:}
\begin{itemize}
\item $r_1$: Defy (resist pressure)
\item $r_2$: Concede (give in)
\item $r_3$: Negotiate (seek compromise)
\item $r_4$: Retaliate (counter-strike)
\item $r_5$: Absorb (accept damage without response)
\end{itemize}

\subsection{Backward Induction Solution}

The Subgame Perfect Equilibrium is found via backward induction:

\begin{enumerate}
\item \textbf{Step 1:} Solve terminal decision nodes (last mover's choice)
\item \textbf{Step 2:} Given Step 1, solve penultimate nodes
\item \textbf{Step 3:} Continue backwards to initial node
\item \textbf{Step 4:} The resulting strategy profile is the SPE
\end{enumerate}


% =============================================================================
% SECTION 5: PARAMETRIZATION VIA LLMMC
% =============================================================================

\section{Parametrization via LLMMC}
\label{sec:crisis-params}

\subsection{LLM Monte Carlo Method}

Parameters are estimated using the LLM Monte Carlo (LLMMC) method:

\begin{enumerate}
\item \textbf{Prior Generation:} LLM generates initial parameter estimates based on:
    \begin{itemize}
    \item Historical analogies (weighted by relevance)
    \item Game theory literature
    \item Domain expertise
    \end{itemize}
\item \textbf{Bayesian Update:} Priors are updated with current evidence:
    \begin{itemize}
    \item News sources (Al Jazeera, Reuters, etc.)
    \item Official statements
    \item Military movements
    \item Economic indicators
    \end{itemize}
\item \textbf{Posterior Output:} Final parameters with uncertainty estimates
\end{enumerate}

\subsection{Parameter Categories}

\begin{table}[h]
\centering
\begin{tabular}{llcc}
\toprule
\textbf{Category} & \textbf{Parameters} & \textbf{Count} & \textbf{Uncertainty} \\
\midrule
Utility Weights & $\alpha_{ij}, \beta_{ik}$ & 12 & $\pm 0.10$ \\
Context Modifiers & $\gamma_k$ & 6 & $\pm 0.15$ \\
Strategic Beliefs & $P(s_j | h)$ & 10 & $\pm 0.15$ \\
Payoff Values & $U_i(z)$ & (derived) & $\pm 0.20$ \\
\bottomrule
\end{tabular}
\caption{Parameter categories in Crisis Escalation Model}
\label{tab:params}
\end{table}


% =============================================================================
% SECTION 6: RESULTS - KEY PREDICTIONS
% =============================================================================

\section{Results: Key Theoretical Predictions}
\label{sec:crisis-results}

The Crisis Escalation Framework yields several testable predictions that distinguish it from traditional international relations models:

\begin{enumerate}
\item \textbf{Context-Dependent Equilibria:} The same game structure produces different equilibria under varying $\Psi$-configurations. High $\Psi_T$ (time pressure) shifts equilibrium toward aggressive moves; high $\Psi_S$ (social observation) moderates face-saving behavior.

\item \textbf{Asymmetric Risk Sensitivity:} Actors exhibit loss aversion ($\lambda > 1$) over status and prestige, leading to escalation even when expected value calculations favor de-escalation.

\item \textbf{Trigger Thresholds:} The framework predicts non-linear responses to triggering events---small changes in observable signals can shift equilibrium paths discontinuously when they cross face-saving thresholds.

\item \textbf{Negotiation Dominance:} Under most parameter configurations, the SPE favors negotiation over conflict, consistent with the observed rarity of major-power wars despite frequent crises.
\end{enumerate}

These predictions are validated through the US-Iran case study in Section \ref{sec:crisis-example}.


% =============================================================================
% SECTION 7: WORKED EXAMPLE: US-IRAN CRISIS (MOD-GEO-001)
% =============================================================================

\section{Worked Example: US-Iran Crisis (MOD-GEO-001)}
\label{sec:crisis-example}

\begin{tcolorbox}[colback=green!5!white, colframe=green!75!black,
    title=Worked Example: US-Iran January 2026]

\textbf{Question:} Wie wahrscheinlich ist ein US-Militärschlag gegen Iran in 48 Stunden?

\textbf{Context (28. Januar 2026):}
\begin{itemize}[nosep]
\item USS Abraham Lincoln + 3 Zerstörer im Mittleren Osten
\item 5'002 Tote bei Protesten im Iran
\item Trump in Davos: «Iran will reden, wir werden reden»
\item Iran: «Finger am Abzug»
\end{itemize}

\textbf{Players:}
\begin{itemize}[nosep]
\item P1: USA (Trump Administration)
\item P2: Iran (Khamenei/IRGC)
\end{itemize}

\textbf{Strategies:}
\begin{itemize}[nosep]
\item USA: \{Threaten, Strike, Negotiate\}
\item Iran: \{Defy, Concede, Negotiate, Retaliate\}
\end{itemize}

\textbf{Context Modifiers:}
\begin{itemize}[nosep]
\item USA: $f(\Psi) = 1.18 \times 0.79 \times 1.14 = 1.06$
\item Iran: $f(\Psi) = 1.38 \times 0.76 \times 0.55 = 0.58$
\end{itemize}

\textbf{Equilibrium (SPE):}
\[
\text{USA: Negotiate} \rightarrow \text{Iran: Accept} \rightarrow \textbf{DEAL}
\]
\[
(U^*_{\text{USA}} = +5.30, \quad U^*_{\text{Iran}} = +2.32)
\]

\textbf{Prediction:}
\[
P(\text{US Strike in 48h}) = 18\% \quad [12\%-25\%]
\]

\textbf{Trigger Scenarios:}
\begin{itemize}[nosep]
\item Iran exekutiert Demonstranten $\Rightarrow P = 45\%$
\item Proxy greift US-Truppen an $\Rightarrow P = 60\%$
\item Direkter Trump-Iran Call $\Rightarrow P = 5\%$
\end{itemize}

\end{tcolorbox}


% =============================================================================
% SECTION 7: INTEGRATION WITH OTHER APPENDICES
% =============================================================================

\section{Integration with Other Appendices}
\label{sec:crisis-integration}

\begin{table}[h]
\centering
\begin{tabular}{lll}
\toprule
\textbf{Appendix} & \textbf{Provides} & \textbf{Used For} \\
\midrule
B (CORE-HOW) & Complementarity $\gamma$ & Interaction structure \\
V (CORE-WHEN) & Context $\Psi$ & Utility modifiers \\
AAA (CORE-WHO) & Actor hierarchy & Player specification \\
C (CORE-WHAT) & Utility dimensions & Benefit/Cost categories \\
BBB (CORE-WHERE) & Parameters $\Theta$ & LLMMC calibration \\
MS (REF-MODELSPACE) & Theory catalog & Game theory basis \\
AN (METHOD-LLMMC) & Estimation method & Parametrization \\
\bottomrule
\end{tabular}
\caption{Integration with EBF Appendices}
\label{tab:integration}
\end{table}


% =============================================================================
% SECTION 8: IMPLICATIONS FOR PRACTICE
% =============================================================================

\section{Implications for Practice}
\label{sec:crisis-practice}

\subsection{For Analysts}

\begin{enumerate}
\item \textbf{Structured Framework:} Use the 9-step model-building workflow
\item \textbf{Context First:} Always specify Ψ-dimensions before parameters
\item \textbf{Uncertainty Quantification:} Report confidence intervals, not point estimates
\item \textbf{Trigger Monitoring:} Track conditional probabilities for key events
\end{enumerate}

\subsection{For Decision-Makers}

\begin{enumerate}
\item \textbf{Equilibrium Insight:} Rational outcome often differs from rhetoric
\item \textbf{Off-Equilibrium Risk:} Triggers can destabilize equilibrium paths
\item \textbf{Context Matters:} Same action, different context = different outcome
\item \textbf{Behavioral Factors:} Impulsivity, face-saving, credibility affect decisions
\end{enumerate}


% =============================================================================
% SECTION 9: SUMMARY
% =============================================================================

\section{Summary}

\begin{tcolorbox}[colback=blue!5!white, colframe=blue!75!black,
    title=Key Takeaways: Crisis Escalation Modeling]

\begin{enumerate}
\item \textbf{Framework:} Sequential Game + Ψ-Context + LLMMC Parametrization
\item \textbf{Solution:} Subgame Perfect Equilibrium via backward induction
\item \textbf{Innovation:} Context modifiers systematically adjust utility functions
\item \textbf{Output:} Probabilities with confidence intervals + trigger scenarios
\item \textbf{Validation:} Checkpoints at predicted timepoints
\end{enumerate}

\textbf{Template Model:} MOD-GEO-001 (US-Iran Crisis)

\textbf{Reusability:} Adaptable to US-China, US-North Korea, Russia-NATO

\end{tcolorbox}


% =============================================================================
% SECTION 10: GLOSSARY
% =============================================================================

\section{Glossary of Symbols}
\label{sec:crisis-glossary}

\begin{tabular}{ll}
\toprule
\textbf{Symbol} & \textbf{Definition} \\
\midrule
$\Gamma$ & Crisis escalation game \\
$N$ & Set of players \\
$H$ & Set of histories \\
$U_i$ & Utility function for player $i$ \\
$U^*_i$ & Context-modified utility \\
$\Psi_k$ & Context dimension $k$ \\
$f_k(\Psi_k)$ & Context modifier function \\
$\alpha_{ij}$ & Benefit weight \\
$\beta_{ik}$ & Cost weight \\
$\gamma_k$ & Context sensitivity parameter \\
SPE & Subgame Perfect Equilibrium \\
LLMMC & LLM Monte Carlo estimation \\
\bottomrule
\end{tabular}

\vspace{5mm}
\textit{For complete symbol definitions, see Appendix G (Central Glossary).}


% =============================================================================
% SECTION 11: CRITICAL FOUNDATIONS
% =============================================================================

\section{Critical Foundations}
\label{sec:crisis-foundations}

This appendix builds on the following foundational elements from the EBF framework:

\begin{enumerate}
\item \textbf{Game Theory (MS-GT-001, MS-GT-002):} Extensive form games with perfect information, solved via backward induction to find Subgame Perfect Equilibrium.

\item \textbf{Context Framework (Appendix UNMAPPED_V):} The eight $\Psi$-dimensions that modify utility functions and beliefs in crisis settings.

\item \textbf{Complementarity (Appendix UNMAPPED_B):} The $\gamma$-matrix capturing how context dimensions interact and reinforce each other in escalation dynamics.

\item \textbf{Parameter Estimation (Appendix UNMAPPED_BBB):} Literature-validated behavioral parameters ($\lambda$, $\sigma$, $\mu$) from Kahneman \& Tversky (1979), Fehr \& Schmidt (1999).

\item \textbf{LLMMC Method (Appendix UNMAPPED_CAL):} LLM Monte Carlo for generating informative priors when empirical data is limited or classified.
\end{enumerate}


% =============================================================================
% SECTION 12: OPEN ISSUES
% =============================================================================

\section{Open Issues and Research Agenda}
\label{sec:crisis-open}

\begin{enumerate}
\item \textbf{Multi-Actor Extension:} Current framework is bilateral; extension to $n > 2$ players with coalition dynamics
\item \textbf{Dynamic Games:} Repeated interactions with reputation and learning
\item \textbf{Information Revelation:} Signaling games with type uncertainty
\item \textbf{Validation Protocol:} Systematic backtesting on historical crises
\item \textbf{Real-Time Updates:} Continuous Bayesian updating with news feeds
\end{enumerate}


% =============================================================================
% SECTION 12: REFERENCES
% =============================================================================

\section{References}
\label{sec:crisis-refs}

\textbf{Core Theory:}
\begin{itemize}
\item Fudenberg, D. \& Tirole, J. (1991). \textit{Game Theory}. MIT Press.
\item Selten, R. (1965). Spieltheoretische Behandlung eines Oligopolmodells. \textit{Zeitschrift für die gesamte Staatswissenschaft}.
\item Schelling, T. (1960). \textit{The Strategy of Conflict}. Harvard University Press.
\end{itemize}

\textbf{Behavioral Extensions:}
\begin{itemize}
\item Kahneman, D. \& Tversky, A. (1979). Prospect Theory. \textit{Econometrica}.
\item Thaler, R. (1980). Toward a Positive Theory of Consumer Choice. \textit{JEBO}.
\end{itemize}

\textbf{International Relations:}
\begin{itemize}
\item Fearon, J. (1995). Rationalist Explanations for War. \textit{IO}.
\item Powell, R. (1999). \textit{In the Shadow of Power}. Princeton.
\end{itemize}

\vspace{5mm}
\textit{For complete bibliography, see Master Bibliography (bcm\_master.bib).}

\nocite{bcm_master}


% =============================================================================
% CROSS-REFERENCE FOOTER
% =============================================================================

\vspace{10mm}
\hrule
\vspace{3mm}

\textbf{Cross-References:}
\begin{itemize}[nosep]
\item $\rightarrow$ Appendix UNMAPPED_B (CORE-HOW): Complementarity framework
\item $\rightarrow$ Appendix UNMAPPED_V (CORE-WHEN): Context dimensions
\item $\rightarrow$ Appendix UNMAPPED_MS (REF-MODELSPACE): Theory catalog
\item $\rightarrow$ Appendix UNMAPPED_LLM (METHOD-LLMMC): Estimation methodology
\item $\rightarrow$ Model Registry: MOD-GEO-001
\item $\rightarrow$ Case Registry: CAS-GEO-001 (pending validation)
\end{itemize}

\vspace{3mm}
\hrule

\end{document}
